%%%%%%%%%%%%%%%%%%%%%%%%%%%%%%%%%%%%%%%%%%%%%%%%%%%%%%%%%%%%%%%%%%%%%
%%%%%%%%%%%%%%%%%%%%%%%%%%%% Anexo %%%%%%%%%%%%%%%%%%%%%%%%%%%%%%%%%%
%%%%%%%%%%%%%%%%%%%%%%%%%%%%%%%%%%%%%%%%%%%%%%%%%%%%%%%%%%%%%%%%%%%%%


% Inicio do Anexo
\begin{anexosenv}

\chapter{História da Linguagem Tex}\label{AxeExe}
\thispagestyle{empty}

Quando o primeiro volume de The Art of Computer Programming de Donald Knuth foi publicado em 1968, o livro tinha sido tipografado com uma composição de metal quente da Monotype Corporation. Este método datado do século XIX tinha um bom estilo clássico, que agrava Knuth. Quando a segunda edição foi publicada em 1976, todo o trabalho precisou ser tipografado novamente porque a tecnologia da Monotype Corporation tinha sido substituída pela fotocomposição e as fontes originais não estavam mais disponíveis. Quando Knuth recebeu as provas da nova versão em 30 de março de 1977, ele as achou terríveis. Naquela época, Knuth viu pela primeira vez a saída de um sistema de tipografia digital de alta qualidade e se interessou pela tipografia digital. O desapontamento com as provas deram a Knuth a motivação final para resolver o problema de uma vez por todas, com o desenvolvimento do seu próprio sistema de tipografia. Em 13 de maio de 1977, ele escreveu um memorando para si mesmo descrevendo as características básicas do TeX.



\end{anexosenv}